\section{Introduction}
\label{sec:introduction}

Multi-protocol DeFi lending allocation has been studied either as a
forecasting problem on hourly-aggregated rate series
\cite{baude2025optimal} or as a control-theoretic problem with
reinforcement learning on the protocol side
\cite{bertucci2025rl, xu2023autogov, agilerate2024}.  Both
formulations share an implicit assumption: that hourly resolution is
sufficient for an allocator to extract economically meaningful edge.
We argue that this assumption is empirically wrong.  Lending-rate
crossovers between Ethereum L1 protocols happen at the speed of pool
deposits and withdrawals --- per-block ($\sim 12$\,s post-PoS), not
per-hour --- because the rate is a deterministic function of
realized utilization, and utilization updates exactly once per
on-chain interaction.  Aggregating to hourly bars destroys the
inter-block microstructure where the allocator's competitive edge
actually lives.

\paragraph{Thesis.}
We argue that DeFi lending allocation is structurally analogous to
high-frequency trading microstructure as described by
\cite{mackenzie2021}, and that the four-class HFT signal taxonomy of
that work transposes onto the multi-protocol lending problem with a
clean mapping.  Concretely: the cross-protocol rate spread is itself
the decision variable (the dominant signal); mempool-visible
transactions are the analog of "order-book dynamics"; lead rates
from Maker DSR and Curve 3pool stand in for the "futures lead" class;
and on-chain gas regime plus stablecoin peg deviations act as the
"related instruments" class.  Under this lens the natural decision
frame is event-time, gas-aware, with three nested policy tiers ---
gas-aware threshold (T1), Ornstein--Uhlenbeck optimal stopping with
switching cost (T2), and Cox proportional-hazards regression on the
signal vector (T3) --- each grounded in a separate microstructure or
quantitative-finance source book
\cite{ohara1995, krause2005, kissell2014, lopezdeprado2018, mackenzie2021}.

\paragraph{Empirical finding.}
On a $\NBlocksFull$ per-block panel covering
\DataStart{}--\DataEnd{} across \NProtocols\ Ethereum USDC lending
venues (\ProtocolList), the gas-aware event-time allocator (T1)
significantly outperforms passive Aave V3 buy-and-hold under a
paired monthly Sharpe bootstrap
($\Delta\textsc{Sharpe} = \HOneAuxDelta$, 95\,\% CI
[\HOneAuxCIlow, \HOneAuxCIhigh], $p = \HOneAuxPvalue$).  On a
\$1\,M position over four months the allocator earns
\$4{,}275 above Aave V3 hold (\$$4{,}039$ above Morpho Blue hold)
at \$$682$ in gas spend over $\TOneNRebal$ rebalances --- two
orders of magnitude more rebalance opportunities than an
hourly-aggregated baseline can express.

\paragraph{Contributions.}
\begin{itemize}
  \item \textbf{An event-time, gas-aware multi-protocol DeFi lending
    allocator with three nested decision-policy tiers (T1/T2/T3),
    each grounded in a distinct microstructure or quantitative-finance
    source.}  We give the closed-form gas-cost crossover inequality
    that locks T1's threshold, the OU--Bellman analytical boundary that
    closes T2, and the Cox proportional-hazards specification with the
    triple-barrier label of \cite{lopezdeprado2018} for T3.
  \item \textbf{A four-class HFT-style signal taxonomy mapped from
    \cite{mackenzie2021} onto multi-protocol DeFi lending} (lead, order
    book, fragmentation, related instruments), with the F3
    (fragmentation) class operationalised as the decision variable
    itself, F1 and F4 as conditioning covariates for T3, and F2
    (mempool) deferred per the explicit risk register.
  \item \textbf{The first per-block-resolution head-to-head matrix of
    seven allocation policies on a multi-month Ethereum mainnet panel}
    (four baselines plus T1/T2/T3), with paired monthly Sharpe
    bootstrap pre-registration, regime-conditional breakdown, and
    Deflated Sharpe Ratio gate per
    \cite[Ch.\,14.7.3]{lopezdeprado2018}.
  \item \textbf{A production-grade reference agent.}  The same
    \texttt{decision/} Python package that produces the empirical
    matrix above is consumed unchanged by an off-chain agent that
    subscribes to \texttt{newHeads}, reads six lending protocols
    per block, and dispatches switches via Flashbots private mempool
    (the asymmetric speed-bump analog of \cite[pp.\,200--203]{mackenzie2021}).
\end{itemize}

\paragraph{Outline.}
\Cref{sec:background} reviews DeFi lending mechanics, the HFT
microstructure literature we draw on, and adjacent work.
\Cref{sec:methodology} specifies the three-tier policy ladder and
the gas-cost crossover inequality.  \Cref{sec:empirical} reports the
seven-policy matrix and the binding $H_{1}^{\text{aux}}$ test plus
the directional secondary tests.  \Cref{sec:discussion} positions
the result against the HFT canon and the cross-domain hinge.
\Cref{sec:limitations,sec:conclusion} state honest scope and next
steps.
